{\color{ink}\bfseries 例 2-12}\par\vspace{1bp}假定初始为空，设\(empty=\allowbreak{}N\)、\(odd=\allowbreak{}0\)、\(even=\allowbreak{}0\)、\(mutex=\allowbreak{}1\)，x为各进程局部变量：\par P1：while(true)\{ x=\allowbreak{}produce();\allowbreak{} P(empty);\allowbreak{} P(mutex);\allowbreak{} put(x);\allowbreak{} V(mutex);\allowbreak{} if(x\%2) V(odd);\allowbreak{} else V(even);\allowbreak{} \}\par P2：while(true)\{ P(odd);\allowbreak{} P(mutex);\allowbreak{} x=\allowbreak{}getodd();\allowbreak{} V(mutex);\allowbreak{} V(empty);\allowbreak{} countodd();\allowbreak{} \}\par P3：while(true)\{ P(even);\allowbreak{} P(mutex);\allowbreak{} x=\allowbreak{}geteven();\allowbreak{} V(mutex);\allowbreak{} V(empty);\allowbreak{} counteven();\allowbreak{} \}\par odd、even分别发放对应产品的取用许可，始终有\(0\le \)实存奇数数+\allowbreak{}实存偶数数\(\le N\)；等待分类许可时未持mutex，生产者可继续放货。题干未写初始存量；若初存o个奇数、e个偶数，则改为\(empty=\allowbreak{}N-o-e\)、\(odd=\allowbreak{}o\)、\(even=\allowbreak{}e\)，且\(0\le o+\allowbreak{}e\le N\)。
